\makeatletter \def\input@path{{../../}} \makeatother
\documentclass[11pt]{article}

\usepackage[utf8]{inputenc}
\usepackage[T1]{fontenc}
\usepackage[spanish,es-nodecimaldot]{babel}
\usepackage[margin=1in]{geometry}
\usepackage{amsmath,amssymb,mathtools,bm}
\usepackage{enumitem}
\usepackage{fancyhdr}
\usepackage[misc]{ifsym}
\usepackage{microtype}
\usepackage{hyperref}

\setlength\parindent{0pt}
\setlength{\headheight}{14pt}
\pagestyle{fancy}
\lhead{\scshape Pontificia Universidad Cat\'olica de Chile}
\rhead{\scshape IMC UC}
\cfoot{\thepage}
\renewcommand{\headrulewidth}{0.1pt}
\renewcommand{\footrulewidth}{0.1pt}
\renewcommand{\labelitemi}{$\blacklozenge$}

\newcommand{\R}{\mathbb{R}}
\newcommand{\p}{\vspace{8pt}}
\DeclareMathOperator{\diver}{div}
\DeclareMathOperator{\Div}{Div}
\DeclareMathOperator{\tr}{tr}
\DeclareMathOperator{\Cof}{Cof}

\begin{document}

\begin{flushleft}
    \small
    \textbf{IMT3130} \hfill
    \textbf{16/05/2026}
\end{flushleft}
\begin{center}
    \LARGE\textbf{Ayudant\'ia 8}\\
    \vspace{3pt}
    \normalsize\textbf{Mec\'anica del continuo, formulaciones mixtas y termodin\'amica}\\
    \normalsize Ayudante: Basti\'an Herrera \Letter{\hspace{1pt}: \texttt{biherrera@uc.cl}}
    \rule{\textwidth}{0.05mm}
\end{center}

\section*{Problemas}

\begin{enumerate}[label=\textbf{\arabic*.}, leftmargin=*]

\item \textbf{Flujo de Stokes--Brinkman incomprensible e inf--sup.}

Sea $\Omega\subset\R^d$, $d\in\{2,3\}$, un dominio Lipschitz acotado y conexo. Se considera un fluido incomprensible, estacionario y de baja velocidad, con velocidad $\bm v:\Omega\to\R^d$ y presi\'on $p:\Omega\to\R$. Sean $\mu>0$ la viscosidad y $\alpha\ge 0$ un coeficiente de arrastre lineal. La fuerza de volumen total se denota por $\bm f\in H^{-1}(\Omega)^d$.

El objetivo es derivar y analizar el modelo
\begin{equation*}
\begin{aligned}
    -\mu\Delta \bm v+\alpha \bm v+\nabla p &= \bm f &&\text{en }\Omega,\\
    \diver \bm v &=0 &&\text{en }\Omega,\\
    \bm v&=\bm 0 &&\text{en }\partial\Omega,\\
    \int_\Omega p\,dx&=0.
\end{aligned}
\end{equation*}

\begin{enumerate}[label=(\alph*), leftmargin=*]
    \item Usando el teorema de transporte de Reynolds, derive la condici\'on local de incomprensibilidad $\diver\bm v=0$ a partir de la conservaci\'on del volumen de todo subcuerpo material.

    \item Partiendo del balance euleriano de momentum lineal
    \begin{equation*}
        \rho\frac{D\bm v}{Dt}=\diver\bm\sigma+\rho\bm b+\bm f_{\mathrm{drag}},
    \end{equation*}
    de la ley constitutiva newtoniana
    \begin{equation*}
        \bm\sigma=-pI+2\mu D(\bm v),
        \qquad
        D(\bm v)=\frac12\left(\nabla\bm v+\nabla\bm v^T\right),
    \end{equation*}
    y del arrastre lineal $\bm f_{\mathrm{drag}}=-\alpha\bm v$, derive el modelo fuerte anterior bajo las hip\'otesis de r\'egimen estacionario lento, viscosidad constante e incomprensibilidad.

    \item Defina espacios adecuados $V$ y $Q$ para la velocidad y la presi\'on. Derive la formulaci\'on d\'ebil mixta y explique por qu\'e la presi\'on se busca con media cero.

    \item Asuma la condici\'on inf--sup continua de la divergencia
    \begin{equation*}
        \inf_{q\in L^2_0(\Omega)\setminus\{0\}}
        \sup_{\bm w\in H^1_0(\Omega)^d\setminus\{\bm 0\}}
        \frac{\int_\Omega q\,\diver\bm w\,dx}
        {\|q\|_{L^2(\Omega)}\|\bm w\|_{H^1(\Omega)}}
        \ge \beta>0.
    \end{equation*}
    Verifique las hip\'otesis necesarias del teorema de Ne\v{c}as--Babu\v{s}ka--Brezzi para concluir existencia, unicidad y una estimaci\'on de estabilidad para $(\bm v,p)$.
\end{enumerate}

\item \textbf{Conducci\'on termoel\'astica con fuente mec\'anica prescrita y Lax--Milgram.}

Sea $\Omega_0\subset\R^d$ un dominio material de referencia. Consideramos la temperatura absoluta $\theta:\Omega_0\times(0,T)\to\R$, con $\theta>0$, y una deformaci\'on prescrita
\begin{equation*}
    \bm\varphi(\bm X,t),
    \qquad
    F=\nabla_{\bm X}\bm\varphi,
    \qquad
    J=\det F>0.
\end{equation*}
La deformaci\'on es conocida, por ejemplo desde un experimento, desde un paso temporal anterior o desde una soluci\'on mec\'anica ya calculada.

Para evitar ambig\"uedades de unidades, escribiremos $R_0$ para el suministro t\'ermico por unidad de volumen de referencia. Se asume que $\rho_0>0$, $c_v>0$, y que el tensor de conductividad $\bm\kappa(\bm X)$ es sim\'etrico, acotado y uniformemente el\'iptico:
\begin{equation*}
    \bm\kappa(\bm X)\xi\cdot\xi\ge \kappa_0|\xi|^2
    \qquad\forall \xi\in\R^d,
    \qquad
    \kappa_0>0.
\end{equation*}
Se impone condici\'on adiab\'atica homog\'enea,
\begin{equation*}
    \bm\kappa\nabla_{\bm X}\theta\cdot\bm N=0
    \qquad\text{en }\partial\Omega_0.
\end{equation*}

\begin{enumerate}[label=(\alph*), leftmargin=*]
    \item Partiendo de la primera ley en configuraci\'on material
    \begin{equation*}
        \rho_0\dot e_m=\bm S:\dot E-\Div_{\bm X}\bm q_0+R_0,
    \end{equation*}
    explique c\'omo se obtiene una ecuaci\'on de calor cuando la potencia mec\'anica se modela como una fuente conocida $g_{\mathrm{mec}}$. Use adem\'as la desigualdad residual de Clausius--Duhem para justificar la ley de Fourier $\bm q_0=-\bm\kappa\nabla_{\bm X}\theta$.

    \item Suponga que el acoplamiento termoel\'astico se eval\'ua expl\'icitamente y produce
    \begin{equation*}
        g_{\mathrm{mec}}
        =3\alpha_TK\,\theta_\star\,F^{-T}:\dot F,
    \end{equation*}
    donde $\alpha_T$, $K$ y $\theta_\star$ son datos conocidos. Use la identidad de la derivada del determinante para escribir $g_{\mathrm{mec}}$ en t\'erminos de $J$ y $\dot J$.

    \item Aplique Euler impl\'icito a
    \begin{equation*}
        \rho_0c_v\dot\theta-
        \Div_{\bm X}\left(\bm\kappa\nabla_{\bm X}\theta\right)
        =R_0+g_{\mathrm{mec}},
    \end{equation*}
    evaluando $g_{\mathrm{mec}}^{n+1}$ como dato conocido. Derive la formulaci\'on d\'ebil para calcular $\theta^{n+1}$.

    \item Pruebe existencia y unicidad de $\theta^{n+1}\in H^1(\Omega_0)$ usando Lax--Milgram. Indique claramente qu\'e hip\'otesis se requieren sobre $R_0^{n+1}$, $g_{\mathrm{mec}}^{n+1}$ y $\theta^n$.
\end{enumerate}

\end{enumerate}

\end{document}
